\chapter{РАЗРАБОТКА АППАРАТНОЙ ЧАСТИ СИСТЕМЫ МОНИТОРИНГА ПАРАМЕТРОВ ОКРУЖАЮЩЕЙ СРЕДЫ}

\section{Разработка структурной схемы устройства}

На основе анализа, проведенного во втором разделе, разработана структурная схема беспроводного узла мониторинга. Основными функциональными блоками являются: модуль управления на базе микроконтроллера ESP32, блок первичных преобразователей (датчики температуры и влажности), блок питания и индикации.

Взаимодействие между микроконтроллером и датчиками осуществляется по цифровой шине I2C, что минимизирует количество используемых проводников и повышает помехоустойчивость системы.

\section{Проектирование принципиальной электрической схемы}

При проектировании схемы особое внимание уделено фильтрации цепей питания. Для подавления высокочастотных помех, возникающих при работе Wi-Fi передатчика, параллельно выводам питания микроконтроллера установлены блокировочные конденсаторы емкостью 0,1 мкФ и 10 мкФ.

В качестве датчика выбран прецизионный сенсор SHT31. Схема его включения предполагает использование подтягивающих резисторов (pull-up) номиналом 4,7 кОм на линиях SDA и SCL. Полная принципиальная электрическая схема узла приведена в приложении к пояснительной записке.

\section{Расчет надежности аппаратного модуля}

Расчет интенсивности отказов устройства $\lambda$ проводится по методике, изложенной в справочных материалах кафедры. Суммарная интенсивность отказов рассчитывается по формуле \eqref{eq:lambda}:

\begin{equation}
    \lambda = \sum_{i=1}^{n} N_i \cdot \lambda_{base, i} \cdot K_E \cdot K_T
    \label{eq:lambda}
\end{equation}

\begin{formula_where}
    $N_i$ --- количество элементов $i$-го типа; \\
    $\lambda_{base, i}$ --- базовая интенсивность отказов элемента; \\
    $K_E$ --- коэффициент эксплуатации; \\
    $K_T$ --- температурный коэффициент.
\end{formula_where}

Наработка на отказ (MTBF) определяется как $T_{ср} = 1 / \lambda$. Для проектируемого устройства расчетное значение $T_{ср}$ составило более 50 000 часов, что соответствует требованиям технического задания.

\section{Конструкторско-технологические решения}

Разработка топологии печатной платы выполнена в системе автоматизированного проектирования KiCad. Плата выполнена по 4-му классу точности. Для обеспечения стабильной работы антенны Wi-Fi в зоне ее расположения предусмотрен вырез в слое металлизации (keep-out zone).



Габаритные размеры устройства в сборе составляют $60 \times 40 \times 20$ мм, что позволяет размещать его в стандартных корпусах для систем промышленной автоматики.

\newpage